% TOP_PROBLEM: {"id": 3070, "title": "The large sets conjecture", "category_id": 2, "category": {"id": 2, "name": "combinatorics", "display_name": "Combinatorics", "description": "Counting problems, graph theory, discrete structures.", "slug": "combinatorics", "order_index": 2, "created_at": "2026-07-31T15:26:25.671Z"}, "status": "open", "problem_number": "OPG-373", "set_id": 12}
% FIRST_POSTED: 2026-10-01
% ATTEMPT_STATUS: unresolved
% UnsolvedMath ID 3070; source catalogue code OPG-373
% !TeX program = lualatex
% GPT-6 Astra Ultra research attempt; independent partial work, 2026-10-01.
\documentclass[11pt,a4paper]{article}
\usepackage[margin=25mm]{geometry}
\usepackage{fontspec}
\setmainfont{lmroman10-regular.otf}[BoldFont=lmroman10-bold.otf,
 ItalicFont=lmroman10-italic.otf,BoldItalicFont=lmroman10-bolditalic.otf]
\usepackage{amsmath,amssymb,mathtools}
\usepackage{unicode-math}
\setmathfont{latinmodern-math.otf}
\AtBeginDocument{\let\setminus\smallsetminus}
\usepackage{xurl,hyperref}
\hypersetup{colorlinks=true,urlcolor=blue}
\setcounter{secnumdepth}{0}
\setlength{\parindent}{0pt}
\setlength{\parskip}{5pt}
\setlength{\emergencystretch}{3em}
\begin{document}
\section*{The large sets conjecture}\label{3070}
{\small UnsolvedMath ID 3070 · OPG-373 · Combinatorics · OpenGarden}
\subsection{Definitions and mathematical statement}
Let $L\subseteq\mathbb N=\{1,2,\ldots\}$. For a positive integer $r$,
an $r$-colouring is any map $\chi:\mathbb N\to\{1,\ldots,r\}$; using
all the colours is not required. A monochromatic arithmetic progression
of length $k\ge2$ with an allowed difference is a sequence
\[
 a,a+d,\ldots,a+(k-1)d,\qquad a\ge1,\quad d\in L,
\]
on which $\chi$ takes a single value. Since $L$ contains positive integers,
these terms are distinct and in increasing order.

The set $L$ is $r$-large if every $r$-colouring has such a progression
for every finite length $k\ge2$. It is large if it is $r$-large for every
positive integer $r$.
The conjecture is
\[
 \forall L\subseteq\mathbb N,\qquad
 L\text{ is }2\text{-large}\ \Longrightarrow\ L\text{ is large}.
\]
Consequently the two-colour hypothesis must force monochromatic
progressions of arbitrary finite length in every finite number of colours,
while retaining the same set $L$ of permitted common differences.

The initial term and difference may depend on both the colouring and the
requested length. Neither a fixed difference that works for every length
nor a single infinite progression is demanded. It is essential that each
individual progression uses one common difference throughout: allowing
successive gaps to vary within $L$ gives a different accessibility problem.
No density, additive closure, decidability, or algebraic structure of $L$
is assumed. A counterexample would be one two-large set failing to be
$r$-large for some finite $r>2$.
\subsection{Short English statement}
If allowed differences force arbitrarily long monochromatic progressions in two colours, must they do so in every finite number of colours?
\subsection{Research attempt: finite witnesses and a periodic obstruction}
\paragraph{Scope and literature check.}
GPT-6 Astra Ultra, 1 October 2026. Allowed differences must remain
inside the same set $L$. The Brown--Graham--Landman problem [S2]
is not a statement about sets with divergent reciprocal sum.
Alweiss [S3] proves that every two-large set is a set of Bohr recurrence,
a substantial necessary recurrence condition. That theorem is credited
as existing literature and is not silently equated with full largeness.
Below we give finite reformulations and an elementary periodic obstruction.

\paragraph{1. Two-largeness has finite certificates at each length.}
For fixed $k$, two-largeness is equivalent to existence of an $N$ such
that every two-colouring of $\{1,\ldots,N\}$ has a monochromatic
$k$-term progression with difference in $L$. The finite-to-infinite
implication follows by restricting a colouring. For the reverse, suppose
no such $N$ exists. Form a tree whose nodes at depth $N$ are the
two-colourings of $\{1,\ldots,N\}$ avoiding the requested progression,
and whose edges are restriction. Every depth is nonempty, and each
node has at most two children. Starting at the root, choose a child
with descendants at arbitrarily large depths; at least one exists
because the branching is finite. Repeating gives an infinite colouring
avoiding the progression, a contradiction.
This argument makes no computability assumption on $L$, and does not
give a uniform bound in $k$ or an algorithm to recognize two-largeness.

\paragraph{2. Every two-large set contains a multiple of each integer.}
Fix $m\ge2$ and factor it as $\prod_{p\mid m}p^{a_p}$. Represent
a residue $x$ by its Chinese-remainder coordinates
$x_p\in\{0,\ldots,p^{a_p}-1\}$. Colour it by
\[
 c(x)=\left(\sum_{p\mid m}
 \left\lfloor x_p/p^{a_p-1}\right\rfloor\right)\bmod2.
\]
No coset of a nontrivial subgroup of $\mathbb Z/m\mathbb Z$ is
monochromatic. To prove this, choose a prime $p$ dividing the subgroup's
order. Its subgroup of order $p$ varies only the $p$-coordinate,
by multiples of $p^{a_p-1}$, while the other coordinates stay fixed.
In such a coset the top $p$-adic digit takes all values $0,\ldots,p-1$,
including both parities. Thus the colouring takes both values there.

Extend $c$ periodically to the positive integers. If a difference $d$
is not divisible by $m$, the residues of a progression of length $m$
with difference $d$ include the entire coset of the nontrivial subgroup
generated by $d$. Such a progression cannot be monochromatic.
Two-largeness applied to this colouring and length $m$ therefore forces
some $d\in L\cap m\mathbb N$. The case $m=1$ just requires
$L\ne\varnothing$, already forced by the definition.

\paragraph{3. Audit of the obvious binary-refinement approach.}
Encode an $r$-colouring using binary digits of its colour labels.
Two-largeness gives an $L$-progression monochromatic in the first
bit, say $a,a+d,\ldots,a+(M-1)d$. To force the next bit constant,
one would apply a colouring argument to the index set. If that
produces indices with common difference $e$, the corresponding
progression in the original integers has difference $de$, not $d$.
Even if $d,e\in L$, the definition gives no reason that $de\in L$.
For example membership of $2$ and $3$ in an arbitrary set does not
force membership of $6$. This example diagnoses a missing inference;
it is not asserted to be a two-large counterexample.

\paragraph{Verification and remaining gap.}
The periodic colouring was checked on every nontrivial subgroup coset
for moduli $2\le m\le40$. The proof supplies the general certificate,
including composite moduli where a simple parity colouring would fail.
Finite compactness permits different thresholds for each requested
length, exactly as the target allows. Neither divisibility recurrence
nor this compactness argument closes the multiplication gap in binary
refinement. A general proof must preserve allowed differences through
colour refinement without adding an unassumed closure property to $L$.

\paragraph{Reproduction of finite checks.}
The finite checks stated above were run with Python 3 using the repository
script [C1]. They verify the stated examples and finite ranges only.

\subsection{Final status}
UNRESOLVED. Finite certificates and a necessary divisibility condition
are proved, and the direct multi-colour induction is stopped at its
precise invalid step. The stronger known recurrence result [S3] is
acknowledged; no assertion that it proves full largeness is made.

\subsection{Sources}
\begin{itemize}
\item[S1] ULAM.AI and the UnsolvedMath Contributors, supplied problems.json, record 3070 (OPG-373), OpenGarden collection. Catalogue identity and source transcription; CC BY 4.0.
\url{https://huggingface.co/datasets/ulamai/UnsolvedMath}.
\item[S2] Open Problem Garden, The large sets conjecture. Original problem and discussion; the mathematical formulation below is expanded from the supplied source record.
\url{http://www.openproblemgarden.org/op/the_large_sets_conjecture}.
\item[S3] R. Alweiss, \emph{$2$-large sets are sets of Bohr recurrence}, arXiv:2512.01997v3, December 2025. Author preprint checked 1 October 2026.
\url{https://arxiv.org/abs/2512.01997}.
\item[C1] Reproducible finite-check code, executed 1 October 2026 with Python 3 (standard library only):
\path{scripts/check_attempt_batch_geometry_2026_10_01.py}.
\end{itemize}
\end{document}
